\subsection{Justificación e Impacto del Proyecto}

La pertinencia del proyecto se sustenta en los beneficios tangibles que aporta a nivel social, económico y organizativo, transformando cualitativamente el servicio de transporte interurbano.

\subsubsection{Justificación Social y para el Pasajero}

En comunidades interurbanas donde la oferta de transporte es discontinua, la certidumbre sobre los desplazamientos impacta directamente en la calidad de vida de las familias:

\begin{itemize}
    \item \textbf{Ahorro de tiempo útil:} Los pasajeros eliminan tiempos muertos prolongados en paradas a cielo abierto, permitiéndoles sincronizar su salida de casa o del trabajo con la aproximación comprobada del micro.
    \item \textbf{Seguridad y bienestar personal:} Reducir la permanencia solitaria en bordes de carretera durante las primeras horas de la mañana o al anochecer previene situaciones de riesgo e inseguridad ciudadana.
    \item \textbf{Equidad en el acceso al servicio:} Garantiza que usuarios de zonas rurales o paradas intermedias tengan la misma oportunidad de solicitar cupo que aquellos ubicados en la parada inicial.
\end{itemize}

\subsubsection{Justificación Operativa y Económica para la Línea y Choferes}

Para los conductores y la organización sindical de la línea, la digitalización se traduce en rendimientos directos:

\begin{itemize}
    \item \textbf{Optimización del gasto de combustible:} Al conocer anticipadamente las paradas solicitadas y evitar aceleraciones o frenadas innecesarias en paradas vacías, el consumo energético de los micros disminuye de forma sistemática.
    \item \textbf{Equidad en la distribución del trabajo:} La asignación cronometrada y transparente de los turnos de salida de 15 minutos evita discordias internas entre choferes sobre turnos preferenciales o adelantamientos indebidos en ruta.
    \item \textbf{Profesionalización institucional:} La secretaría adquiere herramientas formales de auditoría que respaldan la toma de decisiones basada en datos (volumen de viajes, horarios con mayor demanda y tiempos reales de ciclo).
\end{itemize}

\subsubsection{Viabilidad Preliminar del Proyecto}

El proyecto demuestra una alta factibilidad en sus principales dimensiones:

\begin{enumerate}[label=\alph*)]
    \item \textbf{Viabilidad técnica:} El corredor vial Calamuchita -- Tarija cuenta con cobertura sostenida de redes de telefonía móvil de tercera y cuarta generación (3G/4G). Dado que los paquetes de datos para coordenadas GPS y eventos de mano alzada son extremadamente ligeros (en el orden de pocos kilobytes por segundo), la transmisión opera con fluidez aun en condiciones de señal moderada.
    \item \textbf{Viabilidad económica:} Al apoyarse en los dispositivos inteligentes de los propios conductores y usuarios, el proyecto no demanda la adquisición de costosos sensores telemáticos embarcados, manteniendo un costo de implementación accesible.
    \item \textbf{Viabilidad operativa y adopción:} La interfaz simple y la digitalización de una costumbre preexistente (hacer parar el micro con la mano) reducen la curva de aprendizaje, favoreciendo una adopción natural por parte de la comunidad.
\end{enumerate}
